% Page 038

\seite{38}

der Kapelle, sodann Ansprache und\ob
Gesang der Kinder.

\pend
\abschnitt{Revision}
\pstart

Am 11.\ November wurde die hiesige Schule durch\ob
den Herrn Kreisschulinspector revidiert.

\pend
\abschnitt{Revision}
\pstart

Am 20.\ Februar wurde die hiesige Schule durch\ob
den Herrn Kreisschulinspector Wettz\unsicher\ob
revidiert.

\margmark{Seffern, den\\26. März 1904\\Schmitt, Pfarrer.}%
        Osterprüfung.\ob
Am 20.\ März fand durch den Herrn\ob
Ortsschulinspector Schmitt aus Seffern\ob
die übliche Osterprüfung statt.

\pend
\jahresueberschrift{Das Schuljahr 1904}
\pstart

Es wurden entlassen zu Ostern 1904,\ob
4 Knaben. Die Knaben widmen sich der\ob
Landwirthschaft. Aufgenommen wurden\ob
acht Knaben und ein Mädchen. Die\ob
Schülerzahl beträgt nunmehr 60.

\pend
\abschnitt{Vertretung}
\pstart

Auch in diesem Jahre wurde der hiesige Lehrer\ob
zum Kreiskommando nach Trier einberufen. Der\ob
Herr Lehrer in Seffern übernahm die Vertretung.\ob
Der Unterricht fand abwechselnd vormittags\ob
und nachmittags in beiden Schulen statt.

\pend
\abschnitt{Revision}
\pstart

Am 14.\ Juli wurde die Schule durch den Herrn\ob
Kreisschulinspector revidiert.
\pend
